\documentclass[11pt,a4paper]{article}

\usepackage[utf8]{inputenc}
\usepackage[T1]{fontenc}
\usepackage{amsmath,amssymb}
\usepackage{graphicx}
\usepackage{booktabs}
\usepackage{hyperref}
\usepackage[margin=0.9in]{geometry}
\usepackage{caption}
\usepackage{float}
\usepackage{xcolor}

\title{MC Vertex Reweighting for PPG12:\\[2pt]
\large A Data-Driven Iterative Method on the Truth Vertex}
\author{PPG12 Analysis Team (automated report)}
\date{\today}

\begin{document}

\maketitle

\begin{abstract}\noindent
The PPG12 efficiency pipeline must reweight MC to match the measured
$p$+$p$ luminous region, which is very different between the two
crossing angles ($\sigma\!\simeq\!53$\,cm at 0\,mrad vs
$\sigma\!\simeq\!22$\,cm at 1.5\,mrad, Piwinski narrowed). The
production reco-level weight $f(z_r) = D(z_r)/M(z_r)$ closes the reco
vertex distribution exactly, but it leaves the \emph{individual}
truth-vertex widths of double-interaction MC events at a geometric
``half narrow, half wide'' compromise
($\sigma\!\simeq\!47$\,cm instead of 22\,cm), which in turn biases
the truth fiducial cut used by the efficiency denominator. The
resulting cross-section bias at 1.5\,mrad is $+1.4\,\%$ --- comparable
to other PPG12 systematics and well within reach of a correctable
method. This note presents the nominal fix: an \textbf{iterative
data-driven reweight} $w(z)$ applied per truth vertex
($w(z_h)$ for single interactions, $w(z_h)\,w(z_\mathrm{mb})$ for
doubles). The method uses no physics prior for the target shape ---
only the empirical width of the data reco distribution as the seed
Gaussian, and the $99$-th percentile of $|z_r^\mathrm{data}|$ as the
auto-derived closure window. Applied to the photon10 mix, the method
converges cleanly at 0\,mrad (max$|D/M\!-\!1|\!<\!1\,\%$ in 8 iterations)
and reduces $\chi^2$ in the closure window by a factor $\sim\!3100$ at
1.5\,mrad, closing the bulk of the window to the 1--2\,\% level inside
$|z_r|\!<\!30$\,cm and reaching $\sim\!11\,\%$ at the closure-window
edge.
The code is a standalone Python package in
\texttt{efficiencytool/truth\_vertex\_reweight/}; the
per-period \texttt{reweight.root} outputs are ready to be consumed by
\texttt{RecoEffCalculator\_TTreeReader.C}.
\end{abstract}

\tableofcontents

%%========================================================================
\section{Problem and motivation}
\label{sec:problem}
%%========================================================================

PPG12 extracts the isolated prompt-photon cross section from jet-dominated
$p$+$p$ triggered data. The cross section is proportional to the measured
yield divided by the MC-estimated signal efficiency $\epsilon$. Both the
numerator (reco clusters in data, passing the selection) and the
denominator (MC truth photons in the fiducial window) inherit the MC
vertex distribution --- the MC prior on the vertex $z$-coordinate enters
via (i) cluster kinematics reconstructed relative to the reco vertex
$z_r$, (ii) the $|z_r|<\mathtt{vertex\_cut}$ cut applied in data and MC
alike, and (iii) the truth fiducial cut $|z_h|<\mathtt{vertex\_cut}$
used for the efficiency denominator. \emph{Any} data/MC mismatch in the
vertex $z$ distribution (reco or truth) therefore feeds directly into
$\epsilon$ and the extracted cross section.

\paragraph{Notation.}
Throughout this note, $z_r$ denotes the MBD-reconstructed vertex;
$z_h \equiv z_{t,1}$ denotes the primary (hard-scatter) truth vertex;
and $z_\mathrm{mb} \equiv z_{t,2}$ denotes the secondary (minimum-bias)
truth vertex, present only in double-interaction events.
For double-interaction MC, $z_r \approx (z_h + z_\mathrm{mb})/2$ up to
MBD resolution.
Sections~\ref{sec:problem}--\ref{sec:xsec_impact} occasionally use the
$(z_{t,1}, z_{t,2})$ pair notation when the algebraic decomposition into
symmetric/antisymmetric coordinates is needed; elsewhere we use the
$(z_h, z_\mathrm{mb})$ notation exclusively.

\paragraph{The two periods are very different.}
At the 0\,mrad crossing angle the beams overlap along their full length
and the luminous region is $\sigma_z \simeq 53$\,cm. At 1.5\,mrad the
beams cross at an angle, the beam-beam overlap is geometrically shortened
by the Piwinski factor, and the luminous region collapses to
$\sigma_z \simeq 22$\,cm. The MC samples used in PPG12 were generated
with the full $\sim\!65$\,cm beam profile regardless of period; the
reweight has to correct this downward by nearly a factor of three at
1.5\,mrad. This is a very large shape correction and any shortcomings of
the reweight method are amplified at 1.5\,mrad.

\paragraph{The production method.}
The nominal PPG12 reweight
(\texttt{efficiencytool/RecoEffCalculator\_TTreeReader.C} lines
1331--1372) uses a 1D reco-vertex weight
\[
f(z_r) \;=\; \frac{D(z_r)}{M_\mathrm{reco}(z_r)},
\]
where $D(z_r)$ is the data reco-vertex shape and $M_\mathrm{reco}(z_r)$
is the MC reco-vertex shape, both area-normalized. Event-by-event, each
MC event is multiplied by $f(z_r^\mathrm{event})$ using its reconstructed
vertex. By construction,
$M^\prime_\mathrm{reco}(z_r) = f(z_r)\cdot M_\mathrm{reco}(z_r) = D(z_r)$
bin by bin: the reweighted MC reco-vertex distribution is identical to
data.

\paragraph{Why the production method is not enough.}
$f(z_r)$ closes the reco distribution exactly, but it does not close the
individual truth-vertex distributions of double-interaction MC events,
and it is those truth vertices that enter the efficiency denominator via
the truth fiducial cut. The structural reason is a clean algebraic one:
decompose the pair $(z_{t,1},z_{t,2}) = (z_h, z_\mathrm{mb})$ of truth
vertices for a double-interaction event into a symmetric coordinate
$\bar z \equiv (z_{t,1}+z_{t,2})/2$ and an antisymmetric coordinate
$\varepsilon \equiv (z_{t,1}-z_{t,2})/2$. Under the independent Gaussian
sampling of the raw MC, $\bar z$ and $\varepsilon$ are statistically
independent zero-mean Gaussians with
$\sigma_{\bar z} = \sigma_\varepsilon = \sigma_\mathrm{MC}/\sqrt 2$.
The reco vertex is $z_r \approx \bar z$ (up to MBD resolution), so the
production weight $f(z_r)$ depends \emph{only} on $\bar z$ and is flat
in $\varepsilon$.
After reweighting, $\bar z$ matches the target Gaussian
$\mathcal{N}(0,\sigma_\mathrm{lum}^2/2)$ by construction but
$\varepsilon$ is unchanged. The primary truth vertex
$z_h = \bar z + \varepsilon$ is therefore a sum of a narrowed and an
un-narrowed Gaussian,
\begin{equation}
\sigma^2(z_h)
  \;=\; \tfrac{1}{2}\,\sigma_\mathrm{lum}^2 \;+\; \tfrac{1}{2}\,\sigma_\mathrm{MC}^2
  \;=\; \tfrac{1}{2}(22.3^2 + 62.7^2)
  \;\simeq\; (47.0\,\mathrm{cm})^2.
\label{eq:half_narrow_half_wide}
\end{equation}
$f(z_r)$ cannot touch $\varepsilon$, so it cannot narrow $z_h$ and
$z_\mathrm{mb}$ individually below this ``half-narrow, half-wide''
compromise. Any principled fix must therefore act on the truth vertices
directly, i.e.\ it must factorize as
$w(z_h)\,w(z_\mathrm{mb})$ rather than depending only on the reco
combination $z_r$. The nominal method of
Section~\ref{sec:method} does exactly this, with a data-driven
$w(z)$ iterated to closure on the mixed-MC reco distribution.

%%========================================================================
\section{Cross-section impact at 1.5 mrad}
\label{sec:xsec_impact}
%%========================================================================

The efficiency denominator in
\texttt{RecoEffCalculator\_TTreeReader.C} uses both a reco-vertex cut
$|z_r|<\mathtt{vertex\_cut}$ and a truth-vertex cut
$|z_h|<\mathtt{vertex\_cut}$, with
\texttt{vertex\_cut}\,=\,60\,cm in \texttt{config\_bdt\_nom.yaml}. The
reco cut is handled correctly by $f(z_r)$ because the reco distribution
is closed to data exactly. The truth cut is \emph{not} handled correctly
because, as Eq.~\eqref{eq:half_narrow_half_wide} shows, $f(z_r)$ leaves
$\sigma(z_h) \simeq 47$\,cm instead of the target 22\,cm for
double-interaction MC.

\paragraph{Quantitative estimate.} For a Gaussian truth distribution of
width $\sigma$, the fraction in $|z_h|<60$\,cm is
$\mathrm{erf}(60/(\sigma\sqrt{2}))$. Plugging in:
\begin{itemize}\setlength\itemsep{1pt}
\item $\sigma\!=\!47$\,cm (what nominal $f(z_r)$ gives):
      $\mathrm{erf}(60/(47\sqrt 2)) \simeq 0.80$
\item $\sigma\!=\!22$\,cm (correct 1.5 mrad):
      $\mathrm{erf}(60/(22\sqrt 2)) \simeq 0.99$
\end{itemize}
The double-interaction MC loses $\sim\!19$ percentage points of truth
fiducial acceptance that it shouldn't. The mixed MC blend at 1.5\,mrad
is $0.921\!\cdot\!\mathrm{single} + 0.079\!\cdot\!\mathrm{double}$, so
the bias on the integrated efficiency denominator is
\[
0.079 \times 0.19 \;\simeq\; 1.4\,\% \quad \text{at 1.5\,mrad.}
\]
This translates directly to a $+1.4\,\%$ bias on the extracted cross
section at 1.5\,mrad (the efficiency denominator is too small, so the
quoted cross section is too large). At 0\,mrad the analogous bias is
$\sim\!0.3\,\%$ because the MC beam profile is already close to the
(wider) 0\,mrad luminous region. At the percent level the current PPG12
production is biased at 1.5\,mrad because $f(z_r)$ does not correct the
individual truth vertices $z_h$ and $z_\mathrm{mb}$ of double-interaction
events. The bias is comparable to other PPG12 systematics and is easily
fixable with a truth-level reweight.

%%========================================================================
\section{Method: data-driven iterative reweight}
\label{sec:method}
%%========================================================================

The method derives a truth-vertex weight function $w(z)$ --- evaluated at
$z_h$ for single-interaction MC and at both $z_h$ and $z_\mathrm{mb}$
for double-interaction MC --- directly from data, with no Gaussian
assumption on the luminous-region shape, by iterating the mixed-MC reco
distribution to closure against the data reco distribution.

\paragraph{Forward model.}
Let $f_d$ be the cluster-weighted double-interaction fraction in a given
crossing angle (0.079 at 1.5\,mrad, 0.224 at 0\,mrad) and
$f_s = 1 - f_d$ the single-interaction fraction. The mixed MC reco
distribution under a candidate truth-vertex weight $w(z)$ is
\begin{equation}
M(z_r) \;=\;
f_s \!\!\sum_{e \in T_s}\! w(z_h^{e})\,\delta(z_r - z_r^e)
\;+\;
f_d \!\!\sum_{e \in T_d}\! w(z_h^{e})\,w(z_\mathrm{mb}^{e})\,\delta(z_r - z_r^e),
\label{eq:iter_forward}
\end{equation}
evaluated event-by-event over the single-interaction sample $T_s$ and
the double-interaction sample $T_d$. For double-interaction events the
weight enters \emph{twice}, once for the primary truth vertex $z_h$ and
once for the secondary truth vertex $z_\mathrm{mb}$, reflecting the
statistical independence of the two truth vertices under the generator.
This is the crucial structural choice: it gives the iteration access to
the antisymmetric coordinate $\varepsilon$ of
Section~\ref{sec:problem} that the production reco-only weight $f(z_r)$
cannot touch.

\paragraph{Binning and closure window.}
Both periods use the same binning: $80$ bins over $[-200, 200]$\,cm
(bin width 5\,cm), no overflow clipping, and identical bin edges for the
$w(z)$ weight function and the reco-vertex histograms. The iterative closure window is auto-derived per
period as the $99$-th percentile of $|z_r^\mathrm{data}|$:
\begin{itemize}\setlength\itemsep{1pt}
\item 1.5\,mrad: $z_\mathrm{cut}\!=\!83$\,cm (34 bins inside),
      from $|z_r^\mathrm{data}|_{99\%}\!=\!82.0$\,cm.
\item 0\,mrad:   $z_\mathrm{cut}\!=\!144$\,cm (58 bins inside),
      from $|z_r^\mathrm{data}|_{99\%}\!=\!143.9$\,cm.
\end{itemize}
The window is purely data-driven and requires no manual tuning. Outside
it, the statistics in data drop below the level where the iteration
could meaningfully constrain $w(z)$.

\paragraph{Seed.}
The iteration starts from a Gaussian seed
$w_0(z)\!=\!\mathcal{N}(0, \sigma_\mathrm{seed}^2)\,/\,\mathcal{N}(0, \sigma_\mathrm{gen}^2)$,
with $\sigma_\mathrm{gen}$ measured from the photon10 primary truth vertex
$z_h$ (\texttt{vertexz\_truth} branch) ($\sigma_\mathrm{gen}\!=\!64.997 \pm 0.015$\,cm,
intrinsic to the generator and identical between periods, since the
same single MC is loaded for both) and $\sigma_\mathrm{seed} =
\mathrm{std}(z_r^\mathrm{data})$ measured directly from the data
reco-vertex distribution in the relevant run range (21.48\,cm at
1.5\,mrad, 54.32\,cm at 0\,mrad). Note that $\sigma_\mathrm{seed}$
requires \emph{no} physics input --- no Piwinski narrowing factor, no
luminous-region calculation. It is just the empirical width of the
reco distribution we are trying to match.

\paragraph{Iterative update.}
At iteration $i$, the residual is $r_i(z_r) = D(z_r)/M_i(z_r)$, where
$D(z_r)$ is the data reco-vertex histogram and $M_i(z_r)$ is the mixed
MC reco-vertex histogram
from Eq.~\eqref{eq:iter_forward} under the current $w_i$. The update is
Lucy--Richardson-like, multiplicative, and damped:
\begin{equation}
w_{i+1}(z) \;=\; w_i(z)\,\cdot\,\bigl[\mathrm{smooth}(r_i(z))\bigr]^{\alpha},
\qquad \alpha = 0.5,
\label{eq:iter_update}
\end{equation}
where the residual $r_i$ evaluated at reco bin $z_r$ is applied to the
weight at the matching truth bin $z = z_r$ (using the shared binning).
After each update, $w$ is renormalized so that $\langle w(z_h)\rangle$
averaged over the single-MC truth distribution equals unity, preserving
the physics double-interaction fraction.
Damping ($\alpha\!=\!0.5$) prevents over-correction; a 1-bin Gaussian
smoothing on the residual prevents bin-by-bin oscillation; the update
factor is clipped to $[0.5, 2]$ per iteration to bound jumps; and the
endpoints of the update region are edge-clamped to their nearest
interior values to prevent runaway behavior at the boundary. The
update is applied over the full 80-bin binning and the final
$w(z)$ is data-driven everywhere, with no post-iteration snap to a
seed Gaussian. (An earlier version of the method snapped $w(z)$ to
the seed Gaussian outside the closure window; a systematic study with
\texttt{compare\_snap\_strategies.py} showed that removing the snap
improves 0\,mrad closure by a factor $\sim\!1.3$ in $\chi^2$ and
$\sim\!3$ in $\max|r\!-\!1|$ while leaving 1.5\,mrad essentially
unchanged.) Convergence is tested as
$\max|r_i(z)\!-\!1|\!<\!0.01$ inside the closure window, with a hard
cap of 30 iterations.

\paragraph{Implementation.}
The method is a standalone Python tool in
\texttt{efficiencytool/truth\_vertex\_reweight/}. The mixed-MC histogram
is heavy to rebuild ($\sim\!22$\,M events at 1.5\,mrad,
$\sim\!48$\,M at 0\,mrad), so per-sample arrays are cached to
\texttt{output/cache/*.npz} keyed on file mtime; each iteration runs in
seconds per iteration on a single core. A short file list is in
Appendix~\ref{app:code}. Per-period outputs live in
\texttt{output/\{1p5mrad,0mrad\}/} and include
\texttt{reweight.root} (histograms \texttt{h\_w\_iterative},
\texttt{h\_w\_seed}, \texttt{h\_M\_iterative}) and
\texttt{fit\_result.yaml}.

%%========================================================================
\section{Per-period results}
\label{sec:results}
%%========================================================================

Both periods are fit on the photon10 / photon10\_double mix only at
this stage; the full 4-sample blend including jet12/jet12\_double is a
follow-up (Section~\ref{sec:caveats}). Table~\ref{tab:iter_summary}
collects the headline numbers from both fits.

\begin{table}[H]
\centering
\renewcommand{\arraystretch}{1.15}
\begin{tabular}{l c c}
\toprule
quantity & 1.5\,mrad & 0\,mrad \\
\midrule
run range                                & 51274--54000 & 47289--51274 \\
luminosity [pb$^{-1}$]                   & 16.8588      & 32.6574      \\
$f_\mathrm{double}$ (cluster-weighted)   & 0.079        & 0.224        \\
data $|z_r|$ 99\% percentile [cm]        & 82.0         & 143.9        \\
closure window $z_\mathrm{cut}$ [cm]     & 83           & 144          \\
bins inside closure window               & 34           & 58           \\
$\sigma_\mathrm{seed}=\mathrm{std}(z_r^\mathrm{data})$ [cm]  & 21.48 & 54.32 \\
$\sigma_\mathrm{gen}$ from MC $z_h$ (generator) [cm] & \multicolumn{2}{c}{$64.997 \pm 0.015$} \\
\midrule
initial $\chi^2$ (in window)             & $1{,}419{,}454$ & $26{,}589$ \\
final   $\chi^2$ (in window)             & $453$           & $158$ \\
$\chi^2$ reduction factor                & $3132\times$    & $168\times$ \\
final $\chi^2/\mathrm{ndf}$ (in window)  & $13.7$ (ndf 33) & $2.78$ (ndf 57) \\
initial $\max|r-1|$                      & $0.635$         & $0.079$ \\
final   $\max|r-1|$ (in window)          & $0.111$         & $0.008$ \\
iterations used                          & 30 (cap hit)    & 8 \\
converged to $0.01$ tolerance            & \textbf{no}     & \textbf{yes} \\
\midrule
(context) parametric Gaussian $\chi^2/\mathrm{ndf}$ & 3237       & 454 \\
\bottomrule
\end{tabular}
\caption{Iterative data-driven reweight, photon10 mix only, both
crossing angles, with the post-iteration snap to seed removed
(Section~\ref{sec:method}). The 0\,mrad fit converges to the
1\,\% per-bin tolerance in 8 iterations (final
$\max|r\!-\!1|\!\simeq\!0.8\,\%$). The 1.5\,mrad fit reduces $\chi^2$
in the closure window by a factor $\sim\!3100$ in 30 iterations and
stalls at $\max|r\!-\!1|\!\simeq\!11\,\%$, dominated by bins at the
closure-window edge $|z_r|\!\simeq\!83$\,cm. Inside the bulk
$|z_r|\!<\!30$\,cm the per-bin closure is at the 1--2\,\% level
(see Figure~\ref{fig:iter_1p5mrad_closure}). The last row reports the
parametric Gaussian $\chi^2/\mathrm{ndf}$ (best-fit
$\sigma_\mathrm{tgt}$ from Stage-1) against the data; the iterative
method improves on this by a factor of $\sim\!236$ at 1.5\,mrad and
$\sim\!163$ at 0\,mrad.}
\label{tab:iter_summary}
\end{table}

\paragraph{Hero summary.}
Figure~\ref{fig:iter_two_period} is the single-figure summary: data
reco vs reweighted MC reco overlaid for both periods, residuals, the
recovered $w(z)$, and the headline numbers. The 0\,mrad mix-MC matches
data to $\lesssim\!1\,\%$ per bin throughout the closure window; the
1.5\,mrad mix matches to $\sim\!1$--$2\,\%$ in the bulk
$|z_r|\!<\!30$\,cm, degrades to $\sim\!6$--$8\,\%$ around
$|z_r|\!\simeq\!60$\,cm, and reaches $\sim\!11\,\%$ at the
closure-window edge $|z_r|\!\simeq\!83$\,cm.

\begin{figure}[H]
\centering
\includegraphics[width=\linewidth]{figures/vertex_reweight_iter_two_period.pdf}
\caption{Iterative data-driven reweight summary, both crossing angles
($2\times 2$ grid plus summary text). Top-left: data reco $z_r$ (solid)
vs mixed MC reco $z_r$ after iterative reweight (dashed), both periods
overlaid (log scale). Top-right: residual
$r(z_r)\!=\!D(z_r)/M_\mathrm{iter}(z_r)$ inside each period's closure
window (vertical dashed lines). Bottom-left: recovered truth-vertex
weight $w(z)$ for both periods (log scale). Bottom-right: summary
statistics. The 0\,mrad case is essentially closed; the 1.5\,mrad case
is closed in the bulk and stalls at the closure-window edge (see
Section~\ref{sec:interpretation}).}
\label{fig:iter_two_period}
\end{figure}

\paragraph{Truth vertex distribution comparison.}
Figure~\ref{fig:truth_comparison} overlays the primary truth vertex
$z_h$ distribution for the mixed MC under three weighting methods: no
reweight (raw MC, $\sigma\!\simeq\!64$\,cm), the production reco-level
weight $f(z_r)$, and the new iterative truth-level weight $w(z_h)$.
At 1.5\,mrad the production method leaves the mixed-MC primary truth
vertex $z_h$ at $\sigma\!\simeq\!27$\,cm --- narrower than the raw MC
but still far wider than the true luminous-region width
($\sigma_\mathrm{lum}\!\simeq\!22$\,cm), because $f(z_r)$ cannot narrow
$z_h$ and $z_\mathrm{mb}$ individually for double-interaction events
below the half-narrow/half-wide compromise
(Section~\ref{sec:problem}). The iterative method $w(z_h)$ narrows
the $z_h$ distribution to $\sigma\!\simeq\!16$\,cm, which is narrower
than the data reco width $\sigma_\mathrm{seed}\!=\!21.5$\,cm precisely
as expected, since each single-interaction truth vertex must be
narrower than the observed reco width (the MBD vertex finder adds
resolution smearing). At 0\,mrad, where the correction is small, both
methods converge to the same $\sigma\!\simeq\!54$--$57$\,cm.

\begin{figure}[H]
\centering
\includegraphics[width=\linewidth]{figures/vertex_reweight_truth_comparison.pdf}
\caption{Truth vertex $z_h$ distribution for the mixed photon10 MC
under three reweighting methods, per period. Left: 1.5\,mrad, where the
reco-level $f(z_r)$ (blue dashed) fails to narrow the truth
distribution sufficiently, while the iterative $w(z_h)$ (red) narrows
it well past the data reco width (the truth must be narrower than the
reco due to MBD resolution smearing). Right: 0\,mrad, where both
methods agree because the MC beam profile is already close to the data
luminous region.}
\label{fig:truth_comparison}
\end{figure}

\paragraph{0\,mrad: clean convergence.}
Figure~\ref{fig:iter_0mrad_closure} shows the 6-panel diagnostic for
0\,mrad. The seed Gaussian with
$\sigma_\mathrm{seed}\!\simeq\!54$\,cm already captures most of the
shape (initial $\max|r-1|\!=\!8\,\%$, initial
$\chi^2/\mathrm{ndf}\!\simeq\!467$); the parametric best-fit Gaussian
at $\sigma_\mathrm{tgt}\!\simeq\!54.6$\,cm is essentially the same
($\chi^2/\mathrm{ndf}\!\simeq\!454$) because the data is nearly
Gaussian. The iteration cleans up the residual over 8 updates and
drops below the 1\,\% per-bin tolerance inside the closure window
(final $\max|r-1|\!\simeq\!0.8\,\%$, $\chi^2/\mathrm{ndf}\!\simeq\!2.78$).
The recovered $w(z)$ is smooth and broadly Gaussian, with a small
enhancement of the central core relative to the seed, consistent with
the data luminous region being slightly more peaked than a single
Gaussian.

\begin{figure}[H]
\centering
\includegraphics[width=\linewidth]{figures/vertex_reweight_iter_0mrad_closure.pdf}
\caption{0\,mrad iterative closure (6 panels). Top-left/centre: data
reco vs mixed MC reco under no reweight (gray dotted), the seed
Gaussian $w(z)\!\propto\!\mathcal{N}(0,\sigma_\mathrm{seed}^2)/\mathcal{N}(0,\sigma_\mathrm{gen}^2)$
with $\sigma_\mathrm{seed}\!=\!\mathrm{std}(z_r^\mathrm{data})$
(orange dash-dot), the parametric best-fit Gaussian with
$\sigma_\mathrm{tgt}$ from the $\chi^2$ scan (blue dashed), and the
iterative reweight (red). Top-right: the corresponding reweight
functions $w(z)$ (seed, parametric, iterative). Bottom-left: residual
$D/M$ in the closure window with Poisson error bars on the iterative
points and $\pm 1\,\%$/$\pm 2\,\%$ reference bands. Bottom-right:
per-iteration $\chi^2$ and $\max|r-1|$ trajectory. At 0\,mrad the seed
and parametric Gaussians are nearly identical
($\sigma_\mathrm{seed}\!\simeq\!\sigma_\mathrm{tgt}\!\simeq\!54$\,cm)
because the data reco distribution is already close to a single
Gaussian; the iterative method improves on both.}
\label{fig:iter_0mrad_closure}
\end{figure}

\paragraph{1.5\,mrad: bulk closure, edge stall.}
Figure~\ref{fig:iter_1p5mrad_closure} shows the 6-panel diagnostic for
1.5\,mrad. The story is qualitatively different. The seed Gaussian
with $\sigma_\mathrm{seed}\!\simeq\!21$\,cm starts from
$\max|r-1|\!=\!64\,\%$ and an in-window $\chi^2$ of
$1.42\!\times\!10^6$; the first $\sim\!5$ iterations make dramatic
progress, then the iteration enters a slow plateau driven by bins at
the closure-window edge $|z_r|\!\simeq\!83$\,cm.
After 30 iterations, $\chi^2$ in the window has dropped by
$\sim\!3100\times$ to $\simeq\!453$, and $\max|r-1|$ sits at
$\simeq\!0.11$, set by those stuck edge bins. Inside the bulk
$|z_r|\!<\!30$\,cm closure is at the $\sim\!1$--$2\,\%$ level on every
bin; between 30 and 60\,cm the residual climbs to
$\sim\!5$--$10\,\%$. The recovered $w(z)$ shows a strong central core
boost and a clear non-Gaussian shape in the
$|z|\!\simeq\!40$--$60$\,cm region (truth-vertex axis) that no
symmetric Gaussian ansatz can represent.

\begin{figure}[H]
\centering
\includegraphics[width=\linewidth]{figures/vertex_reweight_iter_1p5mrad_closure.pdf}
\caption{1.5\,mrad iterative closure (6 panels), same layout as
Figure~\ref{fig:iter_0mrad_closure}. At 1.5\,mrad the seed Gaussian
($\sigma_\mathrm{seed}\!=\!21.5$\,cm, orange dash-dot) and the
parametric best-fit Gaussian ($\sigma_\mathrm{tgt}\!=\!15.7$\,cm, blue
dashed) are distinct; both are much broader than the data core and
mis-predict the reco-vertex shape by $\sim\!30$--$60\,\%$ at the
closure-window edge. The iterative method (red points) closes the
bulk of the window to $\sim\!1$--$2\,\%$ inside $|z_r|\!<\!30$\,cm and
stalls at $\max|r\!-\!1|\!\simeq\!0.11$ at the closure-window edge
$|z_r|\!\simeq\!83$\,cm. The recovered $w(z)$ in the top-right panel
shows a strong central core boost and a clear non-Gaussian shape in
the $|z|\!\simeq\!40$--$60$\,cm region (truth-vertex axis).}
\label{fig:iter_1p5mrad_closure}
\end{figure}

\paragraph{Method comparison.}
Figure~\ref{fig:iter_1p5mrad_methods} puts the iterative reweight
against three reference baselines at 1.5\,mrad: no reweight, the seed
Gaussian (the $\mathcal{N}(0,\sigma_\mathrm{seed}^2)$ initial point of
the iteration, with $\sigma_\mathrm{seed}\!=\!21.5$\,cm measured from
$\mathrm{std}(z_r^\mathrm{data})$), and the parametric best-fit
Gaussian (a single-parameter fit scanning $\sigma_\mathrm{tgt}$ against
data, best at $\sigma_\mathrm{tgt}\!=\!15.7$\,cm). (Stage 1+2, a
parametric residual correction on top of the parametric Gaussian, is
never used in the nominal pipeline and is not shown.) The no-reweight
MC overshoots data by a factor of several everywhere; both Gaussian
baselines close the width on average but mismatch the shape at the
30--70\,\% level, with visible structure in the ratio panel that no
single-Gaussian ansatz can represent; the iterative method is the only
one that follows the narrow data core while also tracking the
non-Gaussian mid-tail structure.

\begin{figure}[H]
\centering
\includegraphics[width=0.95\linewidth]{figures/vertex_reweight_iter_1p5mrad_methods.pdf}
\caption{Method comparison at 1.5\,mrad: data reco vs mixed MC reco
under no reweight (gray dotted), seed Gaussian
$\sigma_\mathrm{seed}\!=\!21.5$\,cm (orange dash-dot), parametric
best-fit Gaussian $\sigma_\mathrm{tgt}\!=\!15.7$\,cm (blue dashed),
and the iterative data-driven reweight (red), with a ratio panel
below. The two Gaussian baselines bracket the ``best possible'' single-Gaussian
description of the data reco shape; neither closes below the
10--30\,\% level because the true data-MC ratio at 1.5\,mrad is
non-Gaussian. The iterative method is the only one that follows the
narrow data core while also tracking the non-Gaussian mid-tail
structure.}
\label{fig:iter_1p5mrad_methods}
\end{figure}

\paragraph{Parametric $\sigma_\mathrm{tgt}$ scan.}
Figure~\ref{fig:sigma_scan} shows the stage~1 parametric $\chi^2$ scan
that precedes the iterative method. A single Gaussian truth-vertex
weight $w(z) \propto \mathcal{N}(0,\sigma_\mathrm{tgt}^2) /
\mathcal{N}(0,\sigma_\mathrm{gen}^2)$ is evaluated at each scan point.
At 1.5\,mrad the minimum is at $\sigma_\mathrm{tgt}\!\simeq\!15.7$\,cm
with $\chi^2/\mathrm{ndf}\!\simeq\!3237$; the iterative method reaches
$13.7$, a $\sim\!236\times$ improvement. At 0\,mrad the minimum is at
$\sigma_\mathrm{tgt}\!\simeq\!54.6$\,cm with
$\chi^2/\mathrm{ndf}\!\simeq\!454$; the iterative method reaches
$2.78$, a $\sim\!163\times$ improvement. The gap between the best
parametric point and the iterative floor quantifies the shape mismatch
that no single-Gaussian ansatz can capture.

\begin{figure}[H]
\centering
\includegraphics[width=0.95\linewidth]{figures/vertex_reweight_sigma_scan.pdf}
\caption{Stage~1 parametric $\chi^2/\mathrm{ndf}$ as a function of the
target truth-vertex Gaussian width $\sigma_\mathrm{tgt}$, for the
1.5\,mrad (left) and 0\,mrad (right) periods. The minimum of the scan
identifies the best single-Gaussian approximation to the data
reco-vertex shape; the parabolic-refined best-fit
$\sigma_\mathrm{tgt}$ is marked by a vertical dashed line. The
horizontal dashed line (red) shows the final $\chi^2/\mathrm{ndf}$
achieved by the iterative data-driven method, and the dotted line
marks $\chi^2/\mathrm{ndf}=1$ for reference. All $\chi^2$ values are
computed on the reco-vertex axis $z_r$ inside the closure window and
normalized by the number of degrees of freedom.}
\label{fig:sigma_scan}
\end{figure}

%%========================================================================
\section{Physics interpretation}
\label{sec:interpretation}
%%========================================================================

\paragraph{Why 0\,mrad is easy.}
At 0\,mrad the beams overlap head-on along their full length and the
luminous region is broad ($\sigma_\mathrm{lum}\!\simeq\!54$\,cm),
nearly equal to the generator beam profile
($\sigma_\mathrm{gen}\!\simeq\!65$\,cm). The shape of the data reco
vertex is also close to a single Gaussian. The seed is therefore
already a good first approximation, the initial residual is
$\sim\!8\,\%$, and the iteration just cleans up small departures
(slightly enhanced central core). Convergence in a handful of
iterations is what one would expect from a near-textbook deconvolution
problem.

\paragraph{Why 1.5\,mrad is hard.}
Three things conspire at 1.5\,mrad:
\begin{itemize}\setlength\itemsep{2pt}
\item The Piwinski crossing-angle effect compresses the luminous region
to $\sigma_\mathrm{lum}\!\simeq\!21$\,cm, almost a factor of three
narrower than the generator profile. The reweight must extract a
narrow target from a wide source, and the dynamic range of $w(z)$ at
the central core is correspondingly large ($w(0)\!\simeq\!8$).
\item The data reco distribution is \emph{not} a clean Gaussian. It
has a narrow core plus a mid-tail shoulder near
$|z_r|\!\simeq\!45$--$60$\,cm. The iterative method discovers this
structure directly --- the recovered $w(z)$ is non-Gaussian in the
same $z$ range --- whereas any symmetric one-parameter Gaussian ansatz
cannot represent it. This is why the parametric best-fit Gaussian
gives $\chi^2/\mathrm{ndf}\!\simeq\!3237$ on the 1.5\,mrad mix, and
why the iterative method drops it to $\simeq\!14$ (a further
$\sim\!236\times$ improvement).
\item The MBD resolution tail leak from core-truth events fills the
bins near the closure-window edge ($|z_r|\!\simeq\!83$\,cm). The
iteration becomes insensitive to $w(z)$ updates at that $z$ because
the contributions to those reco bins come dominantly from truth vertices
($z_h$, $z_\mathrm{mb}$) at other $z$ values via the wide
double-interaction reco-truth kernel.
The iteration has the wrong handle on those bins, and that is what
sets the residual floor at $\max|r-1|\!\simeq\!11\,\%$.
\end{itemize}

\paragraph{What the recovered $w(z)$ looks like.}
At 1.5\,mrad, $w(z)$ has three distinguishing features visible in the
top-right panel of Figure~\ref{fig:iter_1p5mrad_closure}:
(a) a strong central peak with $w(0)\!\simeq\!8$ that pulls
double-interaction events whose truth vertices ($z_h$ and
$z_\mathrm{mb}$) both lie close to the IP into the data-like core; (b) a mid-$z$ dip that suppresses the
generator's natural $|z|\!\sim\!30$--$50$\,cm population to match the
narrow data core; and (c) a clear $\pm$ asymmetry in
$|z|\!\simeq\!40$--$60$\,cm tracking the asymmetric mid-tail bump in
the data. None of these features is representable by a symmetric
single-parameter Gaussian. At 0\,mrad, $w(z)$ is essentially the seed
Gaussian with a small central enhancement.

%%========================================================================
\section{Caveats and next steps}
\label{sec:caveats}
%%========================================================================

\paragraph{(i) photon10-only mix.}
The current fit uses \texttt{photon10 + photon10\_double} only. The
full production blend also includes \texttt{jet12 + jet12\_double}.
Extending the fit to all four samples is the obvious next step:
mechanically it is two more \texttt{load\_sample} calls plus the
appropriate cluster-weighted blend, but the closure numbers will move
and should be remeasured before any production hookup.

\paragraph{(ii) 1.5\,mrad is not formally converged.}
The 30-iteration cap was hit. Two questions remain open: (a) whether
extending to $\sim\!100$ iterations would push the stuck edge bin to
convergence, or whether it is a hard physical floor set by the
double-interaction reco-truth kernel (the suspicion is the latter);
and (b) whether the failure to converge in $\max|r-1|$ matters for the
downstream analysis. The latter has a fairly clean answer: the PPG12
analysis applies fiducial cuts $|z_r|<60$\,cm on the reco vertex and
$|z_h|<60$\,cm on the primary truth vertex
(Section~\ref{sec:xsec_impact}), and inside those cuts the per-bin
closure is already at the 1--2\,\% level. The stuck edge bin at
$|z_r|\!\simeq\!83$\,cm is well outside the analysis fiducial, so the
iteration is good enough for the analysis use case, even though it is
not formally converged on the full closure-window metric.

\paragraph{(iii) Downstream hookup is wired in but off by default.}
Both \texttt{ShowerShapeCheck.C} and
\texttt{RecoEffCalculator\_TTreeReader.C} now accept two new analysis
fields, \texttt{truth\_vertex\_reweight\_on} and
\texttt{truth\_vertex\_reweight\_file}; when \texttt{on=1}, the code
opens \texttt{reweight.root}, reads \texttt{h\_w\_iterative}, and
multiplies the per-event weight by $w(z_h)$ for single-interaction MC
and $w(z_h)\,w(z_\mathrm{mb})$ for double-interaction MC. The
secondary truth vertex branch \texttt{vertexz\_truth\_mb} (i.e.\ $z_\mathrm{mb}$)
is picked up through an optional-pointer guard, so single-MC trees
without it fall through to the single-vertex path. The two reco-level
and truth-level reweights are mutually exclusive and the config
loader enforces this by forcing \texttt{vertex\_reweight\_on=0}
when \texttt{truth\_vertex\_reweight\_on=1}. The showershape configs
\texttt{config\_showershape\_\{1p5rad,0rad\}.yaml} are pre-wired to
the matching \texttt{output/\{1p5mrad,0mrad\}/reweight.root}; the
nominal BDT configs carry the same fields disabled, so flipping the
main analysis over to the truth-vertex reweight is a one-line config
edit once the jet-sample extension is done.

\paragraph{(iv) Implementation can be sped up.}
Each iteration currently rebuilds the MC reco histogram from scratch
($\mathcal{O}(N_\mathrm{events})$ per pass). Precomputing $z^2$ and
dropping the per-trial mask rebuild would give a $\sim\!10\times$
speedup with no change in result. Not a blocker; nice-to-have for the
all-samples extension.

%%========================================================================
\section{Summary}
\label{sec:summary}
%%========================================================================

$f(z_r)$ cannot close the individual truth-vertex distributions ($z_h$
and $z_\mathrm{mb}$) of double-interaction MC, which biases the primary
truth fiducial cut $|z_h|<60$\,cm used in the efficiency denominator by
$+1.4\,\%$ at 1.5\,mrad
(Section~\ref{sec:xsec_impact}). An iterative data-driven reweight
$w(z)$, applied per truth vertex with the factorization
$w(z_h)\,w(z_\mathrm{mb})$ for double-interaction events, closes this
gap without any Gaussian assumption. Applied to the photon10 mix, the
method converges cleanly at 0\,mrad ($\max|r\!-\!1|\!<\!1\,\%$ in 8
iterations) and reduces in-window $\chi^2$ by a factor $\sim\!3100$
at 1.5\,mrad, closing the analysis-relevant bulk $|z_r|\!<\!30$\,cm to
$1$--$2\,\%$ per bin and reaching $\sim\!11\,\%$ at the closure-window
edge. The recovered $w(z)$ at 1.5\,mrad shows non-Gaussian features
(a strong central boost and a mid-$z$ shoulder) that no closed-form
Gaussian ansatz could represent. Per-period
\texttt{reweight.root} outputs are ready in
\texttt{efficiencytool/truth\_vertex\_reweight/output/} and the C++
hookup into \texttt{ShowerShapeCheck.C} and
\texttt{RecoEffCalculator\_TTreeReader.C} is in place (enabled on the
showershape period configs, wired but off on the nominal BDT
configs). The remaining work is the jet-sample extension of the
standalone fit tool.

\appendix

%%========================================================================
\section{Baseline: events with and without a reconstructed vertex}
\label{app:nonull}
%%========================================================================

Before building any truth-based reweight, we checked that dropping
events whose reco MBD vertex is the $-9999$ sentinel does not itself
bias the MC. The concern is that an MC event which loses its reco
vertex $z_r$ may be physically unusual (e.g.\ with a shifted truth
vertex $z_h$), so dropping it would skew the efficiency in a
$z_h$-dependent way.

\begin{table}[H]
\centering
\begin{tabular}{lrrrr}
\toprule
Sample & $N_\mathrm{total}$ & $N_\mathrm{reco}$ & $N_\mathrm{no\,reco}$ & fraction no reco \\
\midrule
photon10        & 9\,998\,562 & 7\,716\,585 & 2\,281\,977 & \textbf{22.82\,\%} \\
photon10\_double & 9\,997\,075 & 9\,491\,223 &   505\,852 & \textbf{5.06\,\%} \\
jet12           & 10\,000\,000 & 7\,689\,530 & 2\,310\,470 & \textbf{23.10\,\%} \\
jet12\_double   & 9\,997\,000 & 9\,485\,384 &   511\,616 & \textbf{5.12\,\%} \\
\bottomrule
\end{tabular}
\caption{Fraction of events without an MBD reco vertex. The
double-interaction productions have $\sim\!4.5\times$ higher MBD
vertex-finding efficiency than their single-interaction counterparts
(two concurrent collisions give the MBD twice the chance of producing
a vertex).}
\label{tab:counts}
\end{table}

\begin{figure}[H]
\centering
\includegraphics[width=0.95\linewidth]{figures/truth_vertex_reco_check.pdf}
\caption{Primary truth vertex $z_h$ distribution for events with (blue)
and without (red dashed) a reconstructed MBD vertex $z_r$,
area-normalized for shape comparison. Columns: single / double
interaction. Rows: photon10 / jet12. Dotted vertical lines:
$|z_h|=30$\,cm reference.}
\label{fig:shape}
\end{figure}

\begin{figure}[H]
\centering
\includegraphics[width=0.95\linewidth]{figures/truth_vertex_reco_check_ratio.pdf}
\caption{Ratio (no-reco-vertex shape)\,/\,(reco-vertex shape) as a
function of primary truth vertex $z_h$, zoomed to $[0.75,1.25]$. All
samples show a $\sim\!\pm 5\,\%$ north/south asymmetry inside
$|z_h|<30$\,cm, consistent with a small MBD geometric bias.}
\label{fig:ratio}
\end{figure}

Inside $|z_h|<30$\,cm the two shapes agree to within $\pm 5\,\%$ with
a small $z_h$-asymmetry that is the same in all
four samples and therefore a detector-level effect (MBD vertex finder
slightly more efficient at positive $z$). Dropping the no-reco-vertex
events is thus \emph{approximately} harmless --- both the efficiency
numerator and denominator are scaled by the same factor, and the
residual $z$ dependence of the drop rate is small. However, dropping
$\sim\!22\,\%$ of the single-interaction MC sample is still a lot of
statistics, and it makes the efficiency measurement implicitly
conditional on a valid reco vertex being found. That was the original
motivation for exploring a truth-level reweight in the first place.

%%========================================================================
\section{Discarded methods}
\label{app:discarded}
%%========================================================================

Before arriving at the iterative method described in
Section~\ref{sec:method}, we explored three truth-only reweight
constructions that did not work:
\begin{itemize}\setlength\itemsep{1pt}
\item \textbf{$w_1$: direct ratio} of the data reco-vertex shape to the
MC truth-vertex shape, $w_1(z_h) = D(z_r)/T(z_h)$, with the common
binning identifying $z_r$ bins with $z_h$ bins.
\item \textbf{$w_2$: conditional projection},
$w_2(z_h) = \mathbb{E}[f(z_r)\,|\,z_h]$, the expectation of the
production reco weight conditional on the primary truth vertex.
\item \textbf{$w_\mathrm{iter,DA}$: D'Agostini iterative unfolding}
using the MC response matrix $R_{ij}=P(z_r^i|z_h^j)$.
\end{itemize}
All three hit a geometric lower bound on the double-MC sample at
1.5\,mrad: the conditional reco kernel $P(z_r|z_h)$ (for single-interaction MC) has
RMS $\sim\!34$\,cm, wider than the $\sim\!22$\,cm target data peak
(Section~\ref{app:double_diag}), so \emph{any} truth-only reweight is
bounded from below by the kernel width and cannot reach the target.

We also drafted a closed-form \textbf{2D Gaussian reweight}
$w(z_h,z_\mathrm{mb}) \propto \exp[-\tfrac{1}{2}(z_h^2 + z_\mathrm{mb}^2)\,\Delta]$
with $\Delta = 1/\sigma_\mathrm{lum}^2 - 1/\sigma_\mathrm{MC}^2$. This
closed form does work in principle --- it acts on both the symmetric
and antisymmetric coordinates and so can narrow all three vertex axes
simultaneously --- but it assumes a single Gaussian shape for the data
luminous region. At 1.5\,mrad the data has a $\pm$-asymmetric mid-tail
bump that a symmetric Gaussian cannot reproduce, so the closed form
ended up as a first-draft proposal that was superseded by the
iterative method of Section~\ref{sec:method}, which discovers $w(z)$
from data directly without any Gaussian assumption.

%%========================================================================
\section{Double-interaction reco-truth kernel}
\label{app:double_diag}
%%========================================================================

The geometric lower bound on truth-only reweights quoted in
Section~\ref{app:discarded} can be read off the widths directly:

\begin{itemize}\setlength\itemsep{1pt}
\item 1.5\,mrad data reco width: $\sigma\!\simeq\!22.3$\,cm
\item jet12 single-interaction MC truth width ($z_h$): $\sigma\!\simeq\!57.6$\,cm
\item jet12\_double MC truth width ($z_h$, primary): $\sigma\!\simeq\!62.7$\,cm
\item jet12\_double MC reco width ($z_r$):
      $\sigma\!\simeq\!48.6\!\approx\!62.7/\sqrt{2}$ (averaging)
\item Conditional kernel $P(z_r|z_h{=}0)$ RMS (with $z_\mathrm{mb}$
      marginalized): $\sim\!34$\,cm
\end{itemize}

Figure~\ref{fig:double_diag} shows the kernel slice directly,
compared to the 1.5\,mrad data target. Even if the primary truth vertex is forced to
$z_h\!=\!0$, the reco distribution from a double-interaction
event cannot be narrower than the $\sim\!34$\,cm kernel width ---
larger than the 22\,cm target. This is the structural reason the
truth-only methods fail on double-interaction MC, and the structural
reason the iterative method needs the
$w(z_h)\,w(z_\mathrm{mb})$ factorization rather than a
$w$-on-$z_r$ construction.

\begin{figure}[H]
\centering
\includegraphics[width=\linewidth]{figures/vertex_reweight_alt_double_diag.pdf}
\caption{Conditional reco kernel $P(z_r|z_h)$ in jet12\_double MC at
1.5\,mrad (with $z_\mathrm{mb}$ marginalized), with the 1.5\,mrad data
reco $z_r$ shape overlaid. The kernel's
$\sim\!34$\,cm RMS exceeds the $\sim\!22$\,cm data peak, which is the
geometric lower bound on truth-only reweighting.}
\label{fig:double_diag}
\end{figure}

%%========================================================================
\section{Files and code}
\label{app:code}
%%========================================================================

\begin{itemize}\setlength\itemsep{1pt}
\item \texttt{anatreemaker/source/CaloAna24.\{h,cc\}} --- writes
      the \texttt{vertexz\_truth} ($z_h$, primary truth vertex) and
      \texttt{vertexz\_truth\_second} ($z_\mathrm{mb}$, secondary truth
      vertex; read by C++ as \texttt{vertexz\_truth\_mb}) branches that
      the iterative method reads for single- and double-interaction MC
      respectively.
\item \texttt{efficiencytool/truth\_vertex\_reweight/} --- the
      standalone Python tool for the iterative method:
      \begin{itemize}\setlength\itemsep{0pt}
      \item \texttt{fit\_core.py} --- cached sample loader, the
      \texttt{mixed\_mc\_histogram} forward model
      (Eq.~\ref{eq:iter_forward}), Poisson $\chi^2$, and helpers for
      fitting $\sigma_\mathrm{gen}$ from MC truth.
      \item \texttt{fit\_truth\_vertex\_reweight.py} --- the driver
      (stage iterative = Section~\ref{sec:method}), producing
      \texttt{reweight.root} and \texttt{fit\_result.yaml}.
      \item \texttt{plot\_closure.py},
      \texttt{plot\_comparisons.py} --- closure and cross-period
      diagnostic plots.
      \item \texttt{config.yaml} --- sample paths, run ranges, fit
      parameters.
      \end{itemize}
\item \texttt{efficiencytool/TruthVertexRecoCheck.C} ---
      Appendix~\ref{app:nonull} supporting macro.
\item Outputs:
      \texttt{efficiencytool/truth\_vertex\_reweight/output/\{1p5mrad,0mrad\}/reweight.root},
      \texttt{fit\_result.yaml}, and the cached event arrays in
      \texttt{output/cache/}.
\end{itemize}

\end{document}
